\documentclass[11pt]{article}
\usepackage{amsmath,amssymb,amsthm,booktabs,graphicx,natbib,hyperref,geometry}
\geometry{margin=1in}
\newtheorem{proposition}{Proposition}

\title{Strategic Interaction among Public Innovation Hubs under Private Repricing}
\author{}
\date{}
\begin{document}
\maketitle
\begin{abstract}
\textbf{Introduction.} We study two regional governments that invest in public innovation intermediation while sharing a price-setting metropolitan private alternative. Cross-side participation can make public investments strategic complements when the private access price is held fixed, whereas endogenous private repricing adds an opposing strategic channel. \textbf{Methods.} We compare a price-matched fixed-price benchmark with the full sequential game, derive the public cross effects analytically around the zero-network benchmark, and independently audit a repaired baseline witness over the full public strategy interval and all project-routing regimes. \textbf{Results.} On regular interior stationary branches, the fixed-price public cross effect is zero at $\beta=0$ but has a strictly positive first derivative in $\beta$, while the full-game cross effect is already strictly negative at $\beta=0$. Hence there exists $\varepsilon>0$ such that, for sufficiently small positive network interaction on those continuing branches, the fixed-price local best-response slope is positive while the full-game local slope is negative. Separately, an all-regime computational audit at a repaired baseline vector finds a global public best response in each environment and reproduces the opposite local slopes. \textbf{Conclusion.} Endogenous commercial repricing can change the sign of decentralized public strategic interaction relative to a price-matched benchmark. The analytic result is local and sufficient rather than a global primitive-space theorem; the repaired global-equilibrium result is a computational witness rather than a general theorem.
\end{abstract}
\noindent\textbf{Keywords:} public investment; two-sided platforms; strategic interaction; local public economics; endogenous pricing; innovation intermediation; network effects.\\
\noindent\textbf{JEL Classification:} H77; L13; D62; O38.

\section{Introduction}\label{sec:introduction}

Regional governments increasingly use public organizations to connect firms, universities, investors, experts, and other partners around innovation projects. At the same time, firms outside major cities can access dense commercial innovation ecosystems in metropolitan areas. These two forms of intermediation need not be pure competitors at the ecosystem level: they may collaborate, exchange participants, and provide complementary services. Yet for a focal project deciding which intermediary will be its primary route to partners, the alternatives can be substitutes. This coexistence raises a strategic question for decentralized public investment. When one region strengthens its own intermediation capacity, does that induce another region to invest more or less?

The answer is not immediate because innovation intermediation is two-sided. A regional investment that attracts or facilitates partners raises the value of the public route to projects. More projects then make the same route more attractive to additional partners. With a shared private alternative, an investment in one region also shifts projects and partners away from the private route. If the private access price is fixed, that weakening of the shared outside option can raise the marginal return to public investment in the other region. The public investments can therefore be strategic complements.

The private alternative, however, need not passively keep its price unchanged. A commercial innovation intermediary can respond to lost demand by cutting its access price. This response partly restores the private route's attractiveness to projects in both regions. The induced repricing channel works against the positive cross-side network channel. The question of this paper is whether it can be strong enough to change the qualitative strategic relationship between the two decentralized public investors.

We develop a two-region model with two public innovation intermediaries and a shared metropolitan private intermediary. Regional governments simultaneously choose facilitation investments. The private intermediary then chooses a project-side access price. Projects select one primary intermediation route, while support-side partners may participate in multiple ecosystems. Project participation and partner participation reinforce one another. The governments maximize regional welfare; the private intermediary maximizes access-fee profit.

Our main analytic result is a local strategic sign reversal. We compare the full game with a nested benchmark that holds the private price fixed at the full-game equilibrium fee. Consider regular interior stationary branches around the zero-network benchmark, with the G branch anchored at a symmetric regular beta-zero full-game stationary state. In the fixed-price benchmark B3, the public cross derivative is exactly zero at $\beta=0$, but its first derivative with respect to $\beta$ is strictly positive. In the full game G, by contrast, the public cross derivative is already strictly negative at that symmetric regular state when $\beta=0$. Continuity and strict second-order conditions therefore imply that there exists $\varepsilon>0$ such that, for every sufficiently small positive $\beta$ on the continuing regular stationary branches,
\[
\BR_i^{B3\prime}>0>\BR_i^{G\prime}.
\]
This is a sufficient local theorem about the strategic geometry of continuing interior stationary branches. It is not a theorem that those branches are globally optimal for every nearby primitive vector, it does not characterize the entire primitive parameter region, and it does not provide an explicit formula for $\varepsilon$.

The economic reason for the reversal is transparent. With the private price fixed, cross-side participation transmits one region's public expansion through the shared outside option and raises the other region's marginal return to facilitation. Once the commercial intermediary is allowed to optimize its price, public expansion also changes private demand and therefore the private pricing policy. The resulting price response feeds back into project routing in both regions. In derivative form, the full-game public cross effect equals the fixed-price effect plus terms generated by the private price policy and its first and second derivatives. When those repricing terms dominate the fixed-price network force, the public game changes locally from strategic complementarity to strategic substitution. The negative follower-price contribution is not itself claimed as novel; the contribution is the qualitative reversal in the same public--public strategic object when the private pricing policy is activated.

We complement the analytic theorem with a separate repaired global-equilibrium witness. The earlier rational vector used in the first manuscript draft admits profitable finite deviations once all routing regimes are solved and therefore is not used as global-equilibrium evidence. The repaired baseline keeps the model architecture unchanged and uses $\beta=0.01$, $\gamma=0.825$, and $\tau=0.35$, with the other primitives unchanged. An implementation-independent all-regime audit reoptimizes the private continuation after public deviations, searches the full public strategy interval, and checks the previously dangerous low-investment and boundary regions from multiple participation starts. It finds symmetric public best responses near $x_G=0.83710$ in G and $x_{B3}=0.82589$ in B3 and reproduces opposite local slopes, approximately $-0.0199$ and $+0.0040$, respectively. This is computational certification of one repaired baseline witness, not a primitive-space theorem or a uniqueness theorem over parameter values.

The matched-price construction is central to interpretation. B3 uses the same on-path private fee chosen in G but suppresses the private price policy response to public deviations. The comparison therefore isolates the repricing-response margin rather than a difference in the level of the private fee. B3 is an identification benchmark, not a literal price-regulation counterfactual.

Welfare is deliberately separated from strategic interaction. The private access fee is a regional project outflow but an aggregate transfer to the private intermediary, so changing the fee does not mechanically change national welfare. At the repaired full-game witness, a marginal increase in one region's public investment creates a positive effect on the other region's reduced welfare that is only partly offset by lower private profit. The resulting local national coordination wedge is positive. This is a baseline-specific marginal under-provision result; it is not a solved first-best comparison and does not imply that strategic substitution is socially desirable.

The institutional interpretation is also intentionally narrow. Public innovation hubs and commercial metropolitan innovation communities in Japan and abroad exhibit the relevant ingredients: facilitation and matching, cross-side communities, paid access for some private organizations, and access by projects outside the metropolitan core. These examples support the plausibility of the primitives but do not establish the model's causal repricing prediction. In particular, project single-homing is a focal-project primary-route abstraction, private profit maximization is a reduced-form representation of a selected commercial operator class, and the repaired value of the remote-public access friction is constructive rather than empirically calibrated.

Our contribution is not the discovery of network complementarity, endogenous follower pricing, mixed public/private competition, public-investment strategic substitution, or two-sided platforms. Each of those components has close precedents. Sequential-investment and fiscal-competition literatures already show that downstream strategic responses can alter upstream investment incentives, while platform work studies pricing, network feedback, and public/private competition. The narrower contribution is the matched-price comparison for the same two decentralized public investment choices: the local public--public slope is positive when the shared private price response is switched off at the matched full-game fee and negative when that private price setter is activated. The local theorem and the repaired all-regime computational witness support that proposition without converting the familiar component mechanisms into novelty claims.

The rest of the paper proceeds as follows. Section \ref{sec:model} presents the model. Section \ref{sec:equilibrium} characterizes participation, pricing, and the nested benchmarks. Section \ref{sec:mainresult} develops the mechanism decomposition, states the analytic local sign-reversal theorem, and reports the repaired computational witness. Section \ref{sec:welfare} develops welfare and coordination. Section \ref{sec:robustness} separates proved local scope from unproved extension directions. Section \ref{sec:institution} discusses institutional interpretation and empirical implications. Section \ref{sec:literature} positions the result in the related literature. Section \ref{sec:discussion} discusses limitations and scope, and Section \ref{sec:conclusion} concludes.

\section{Model}
% Stage 10 substantive writing only.

\section{Equilibrium and Benchmarks}\label{sec:equilibrium}

This section reduces the participation subgame to a smooth system and defines the two nested environments used to identify the role of private repricing.

\subsection{Participation fixed point}

Let
\[
\delta=\alpha\beta,\quad
A_i=v+\alpha(\rho+x_i),\quad
A_T=v+\alpha\rho_T,\quad
d=\kappa_L-\kappa_T-p_T .
\]
In the central regime, local project mass is $n_i^F=1-t_i$ and private project mass is
\[
n_T^F=t_1+t_2-\frac{2(\kappa_T+p_T)}{q_T}.
\]
Using \eqref{eq:partnerlocal}--\eqref{eq:partnerprivate}, the project-quality indices satisfy
\begin{align}
q_i &= A_i+\delta(1-t_i), \label{eq:qi}\\
q_T &= A_T+\delta\left(t_1+t_2-\frac{2(\kappa_T+p_T)}{q_T}\right). \label{eq:qT}
\end{align}
Together with the indifference condition $t_i(q_i-q_T)=d$, the participation equilibrium solves
\begin{align}
F_i &\equiv t_i\left[A_i+\delta(1-t_i)-q_T\right]-d=0,\qquad i=1,2, \label{eq:Fi}\\
F_T &\equiv q_T-A_T-\delta\left[t_1+t_2-\frac{2(\kappa_T+p_T)}{q_T}\right]=0. \label{eq:FT}
\end{align}
This three-equation system is the smooth object used for implicit differentiation and numerical verification within the central interior regime.

\subsection{Private pricing}

Let $D_T\equiv n_T^F$ denote private-route project demand after participation has adjusted. An interior private price satisfies
\begin{equation}
D_T+p_TD_{T,p}=0,
\label{eq:privatefoc}
\end{equation}
with second-order condition
\begin{equation}
2D_{T,p}+p_TD_{T,pp}<0.
\label{eq:privatesoc}
\end{equation}
Implicit differentiation of \eqref{eq:privatefoc} yields the policy response
\begin{equation}
p^*_{T,x_i}
=
-\frac{D_{T,x_i}+p_TD_{T,px_i}}
{2D_{T,p}+p_TD_{T,pp}}.
\label{eq:px}
\end{equation}
At the reported full-game stationary candidate this derivative is negative. Economically, stronger regional facilitation draws projects away from the shared private route; the private intermediary responds by lowering its optimal access price.

The negative price response is not itself a new mechanism. Indeed, when $\beta=0$ the private demand problem becomes linear and the sign of the price response can be established directly. The role of two-sidedness in our main result is different: it creates the positive public--public cross effect in the fixed-price benchmark.

\subsection{Nested benchmarks}

We compare two environments with identical primitives.

\paragraph{Fixed-price benchmark (B3).}
Both regional public hubs and the shared private hub operate, and participation remains two-sided, but the private access price is held fixed while a public deviation is evaluated. Along the G stationary branch used by the local theorem, define the matched scalar fee
\begin{equation}
\bar p_T(\beta)
= p_T^*\!\left(x_1^G(\beta),x_2^G(\beta);\beta\right).
\label{eq:matchedpricepath}
\end{equation}
For each fixed $\beta$, every B3 public deviation is evaluated holding this scalar $\bar p_T(\beta)$ fixed. The public choices and resulting allocation are otherwise allowed to readjust within B3. Hence B3 and G use the same matched scalar fee for the comparison, but their stationary public investments generally differ: $x_i^{B3}(\beta)$ need not equal $x_i^G(\beta)$.

This distinction also separates two derivative objects. The chain-rule decomposition of the full-game objective can switch the price-policy derivative off locally at a given G state. The headline B3--G comparison instead compares the local strategic slopes at the respective B3 and G stationary states associated with the same matched fee path. Those are not the same-state allocation.

B3 is therefore an identification benchmark for the private repricing margin. It suppresses the policy response $p_T^*(x_1,x_2)$ while avoiding a mechanical difference in the matched fee level. It is not a planner problem, a literal price-regulation counterfactual, or a claim that public investments are complements for every exogenous value of $\bar p_T$.

\paragraph{Full game (G).}
The private intermediary observes $(x_1,x_2)$ and optimally chooses $p_T$ according to \eqref{eq:privatefoc}. Governments anticipate the reduced price policy $p_T^*(x_1,x_2)$ when choosing facilitation investments.

Let $\widetilde W_i$ denote government $i$'s reduced objective after downstream private pricing and participation have been solved. At an interior local optimum,
\begin{equation}
\BR_i'(x_j)
=
-\frac{\widetilde W_{i,x_ix_j}}
{\widetilde W_{i,x_ix_i}}.
\label{eq:brslope}
\end{equation}
Because the own second derivative is negative at the regular interior configurations considered below, the sign of the best-response slope is the sign of the public cross derivative.

\section{Strategic Interaction and Private Repricing}\label{sec:mainresult}

This section separates three objects that play different roles in the argument: the fixed-price network force, the additional repricing force, and the resulting sign comparison between B3 and G. The first two explain the mechanism; the third is the main analytic result. A separate repaired computational witness establishes that the same qualitative comparison is compatible with global public best responses in the baseline model.

\subsection{Fixed private price: cross-side reinforcement}

Let
\[
M_{B3}\equiv W_{i,x_ix_j}^{B3}
\]
denote the public cross derivative after participation has been solved while the private access price is held fixed. Under a strict public second-order condition, \eqref{eq:brslope} implies that $M_{B3}$ and the local B3 best-response slope have the same sign.

When $p_T$ is fixed, an increase in $x_j$ raises support-side participation at $H_j$, attracts projects toward the regional public route, and reduces project and partner participation at the shared private route. The weakened shared outside option then raises the marginal value of public facilitation in region $i$. Cross-side participation is the only cross-region link in this benchmark central regime.

\begin{proposition}[First-order benchmark complementarity]\label{prop:b3local}
Consider a regular interior B3 stationary branch through $\beta=0$ with $d=\kappa_L-\kappa_T-\bar p_T>0$ and beta-zero quality gaps $\Delta_i=A_i-A_T>0$. Then
\[
M_{B3}(0)=0
\]
and
\begin{equation}
\left.\frac{\partial M_{B3}}{\partial\beta}\right|_{\beta=0}
=
\frac{\alpha^3d^3}{\Delta_i^3\Delta_j^2}>0.
\label{eq:b3betaderivative}
\end{equation}
Consequently, if the stationary branch and its strict public second-order condition continue smoothly, the local B3 best-response slope is positive for all sufficiently small positive $\beta$ on that branch.
\end{proposition}

The result identifies exactly what two-sidedness contributes. At $\beta=0$, B3 has no public--public cross effect. A positive network interaction turns that zero cross effect positive to first order. Because $M_{B3}(0,x,p)$ is identically zero within the beta-zero central regime, derivatives of the stationary path and of the matched price do not enter the first derivative in \eqref{eq:b3betaderivative}; the positive term is the direct first-order cross-side effect.

\subsection{Endogenous private price: the repricing channel}

In the full game, government $i$ evaluates
\[
\widetilde W_i(x_i,x_j)
=
W_i(x_i,x_j,p_T^*(x_i,x_j)).
\]
Write $p_i=p^*_{T,x_i}$, $p_j=p^*_{T,x_j}$, and $p_{ij}=p^*_{T,x_ix_j}$. Exact differentiation gives the decomposition
\begin{equation}
M_G
=
M_{\mathrm{fixed}}
+
P_{\mathrm{price}},
\qquad
P_{\mathrm{price}}
=
W_{i,x_ip}p_j+W_{i,x_jp}p_i+W_{i,pp}p_ip_j+W_{i,p}p_{ij},
\label{eq:repricingdecomposition}
\end{equation}
where $M_{\mathrm{fixed}}=W_{i,x_ix_j}$ is the cross derivative that would remain if the price policy were locally switched off at the same state.

\begin{proposition}[Private-price feedback]\label{prop:pricefeedback}
At any regular full-game state with $M_{\mathrm{fixed}}>0$, endogenous repricing reverses the sign of the public cross derivative if and only if
\[
-P_{\mathrm{price}}>M_{\mathrm{fixed}}>0.
\]
Thus repricing changes local strategic complementarity into local strategic substitution precisely when the price-policy contribution dominates the positive fixed-price force at that state.
\end{proposition}

Proposition \ref{prop:pricefeedback} is an identity-based mechanism statement rather than a novelty claim about follower pricing. Economically, the chain is
\[
\begin{aligned}
x_j\uparrow
&\Longrightarrow \text{private demand falls}\\
&\Longrightarrow p_T^*\text{ adjusts}\\
&\Longrightarrow \text{project routing in region }i\text{ changes}\\
&\Longrightarrow W_{i,x_i}\text{ changes}.
\end{aligned}
\]
At the repaired baseline witness, the model predicts $p^*_{T,x_i}<0$: stronger regional facilitation induces the private intermediary to cut its price. This price cut restores part of the private route's attractiveness and works against the positive fixed-price network force.

\subsection{Analytic local strategic reversal}

The beta-zero limit makes the contrast between the two forces especially sharp. In B3 the cross effect starts from zero and moves positively with network interaction. In G, private repricing already couples the governments when $\beta=0$.

\begin{theorem}[Local strategic sign reversal]\label{thm:reversal}
Suppose there exist regular interior B3 and G stationary branches at $\beta=0$. Assume the relevant participation and stationary-condition Jacobians are nonsingular, the private stationary price and both public stationary choices satisfy strict second-order conditions, and the branches continue smoothly for small $\beta$. Let the B3 benchmark use the matched full-game price. Then there exists $\varepsilon>0$ such that, for every $\beta\in(0,\varepsilon)$ on these continuing stationary branches,
\[
\BR_i^{B3\prime}>0>\BR_i^{G\prime}.
\]
More specifically, B3 satisfies \eqref{eq:b3betaderivative}, while at a symmetric regular beta-zero full-game stationary state, with $T=A_T=v+\alpha\rho_T>0$ and $\Delta=A_i-A_T>0$,
\begin{equation}
M_G(0)
=
-\frac{3T\alpha^2\kappa_L^2}
{16\Delta(\Delta+T)^3}<0.
\label{eq:gzero}
\end{equation}
\end{theorem}

The proof is in Appendix \ref{app:smallbeta}. The theorem is a sufficient local result about continuing interior stationary branches. It does not by itself prove global public optimality along those branches, provide a primitive formula for $\varepsilon$, establish uniqueness, or characterize the entire parameter space.

\begin{corollary}[Local strategic effect of commercial adaptation]\label{cor:interpretation}
On the regular branches covered by Theorem \ref{thm:reversal}, holding the private alternative's price fixed can give the wrong qualitative description of the local public strategic relationship. Activating the private intermediary's optimizing price policy can change not only the magnitude but also the sign of the local public--public best-response slope.
\end{corollary}

\subsection{Repaired all-regime computational witness}

The analytic theorem is not used as a global-equilibrium proof. Global best-response status is checked separately at a repaired baseline vector that keeps the model architecture unchanged. The repaired witness uses
\[
\beta=0.01,\qquad \gamma=0.825,\qquad \tau=0.35,
\]
with the remaining primitives as reported in Appendix \ref{app:tables}. At this vector, an implementation-independent evaluator reconstructs all four project routes from the utility upper envelope, solves support-side participation from multiple starts, reoptimizes the private price after public deviations, searches the full public strategy interval $[0,1]$, and intensively rechecks boundaries and the low-investment region that invalidated the earlier draft witness.

The audit finds symmetric public best responses near
\[
x_G\simeq0.837102,\qquad x_{B3}\simeq0.825890,
\]
with a matched private price
\[
p_G\simeq0.018404.
\]
The best detected deviation gain is below the documented numerical certification tolerance in both environments. Local derivative checks at the same repaired witness give approximately
\[
\BR_i^{G\prime}\simeq-0.0199,
\qquad
\BR_i^{B3\prime}\simeq+0.0040.
\]
Thus the baseline model contains a computationally certified all-regime witness with the same qualitative sign reversal. This is an existence witness under the baseline functional form and documented numerical protocol. It is not an exact interval proof of global optimality for all primitives, a uniqueness theorem over parameter values, or a general robustness theorem.

The remote-public access condition at the repaired witness is intentionally transparent:
\[
\kappa_L+\tau>v+\alpha.
\]
It makes a nonresident rival public hub dominated by non-participation for every project type and removes the direct free-riding route that destroyed the earlier draft witness. The value of $\tau$ is constructive rather than empirically calibrated. Consequently, the repaired global witness should not be interpreted as evidence of direct user poaching between the two public hubs; the certified public--public strategic interaction operates through the shared private route and cross-side participation.

Table \ref{tab:strategic} reports the repaired witness magnitudes from the v2.1 result pipeline.

\begin{table}[t]
\centering
\caption{Strategic interaction at the repaired baseline witness}
\label{tab:strategic}
% DO NOT EDIT — GENERATED FILE
\begin{tabular}{lrr}\toprule Environment & reported $x$ & local BR slope \\ \midrule
G & 0.837102 & -0.019874 \\
B3 & 0.825890 & 0.003968 \\
\bottomrule\end{tabular}

\end{table}

An earlier exact rational Krawczyk calculation remains archived in the repository as a local stationary-root diagnostic for the rejected draft vector. Because that draft vector admits profitable finite deviations once omitted routing regimes are restored, the exact stationary-root calculation is not used as global-equilibrium evidence in this paper.

\section{Welfare and Coordination}\label{sec:welfare}

The sign of a government's best-response slope is not a welfare ranking. Strategic interaction describes how one government's optimal investment responds locally to the other's choice; welfare depends on project allocation, partner surplus, real facilitation costs, and private profit. We therefore keep the strategic reversal and the welfare analysis separate.

\subsection{Regional and aggregate welfare}

In the central regime, project surplus in region $i$ is
\begin{equation}
PS_i
=
\int_s^{t_i}(zq_T-\kappa_T-p_T)\,dz
+
\int_{t_i}^{1}(zq_i-\kappa_L)\,dz .
\label{eq:PS}
\end{equation}
With uniform support-side costs, the total surplus associated with partner participation at route $h$ is $(n_h^P)^2/2$. Under the symmetric regional attribution used in the regional objectives, each region receives one half of that national partner surplus. Regional welfare is therefore
\begin{equation}
W_i
=
PS_i
+
\frac{(n_1^P)^2+(n_2^P)^2+(n_T^P)^2}{4}
-
\frac{\gamma}{2}x_i^2.
\label{eq:Wi}
\end{equation}
Aggregate welfare is
\begin{equation}
W^N=W_1+W_2+\Pi_T.
\label{eq:WN}
\end{equation}

A project using $H_T$ pays $p_T$ in regional project surplus, while the same payment appears as private profit. Because private operating cost is normalized to zero, the transfer cancels exactly at the aggregate level:
\[
-p_Tn_T^F+\Pi_T=0.
\]
Accordingly, a lower private price does not mechanically imply higher aggregate welfare.

\subsection{A local cross-regional coordination wedge}

At the repaired interior decentralized full-game witness, government $i$'s own reduced first-order condition is numerically zero within the documented tolerance. The national marginal effect of increasing $x_i$ can therefore be decomposed as
\begin{equation}
\frac{dW^N}{dx_i}
=
\frac{d\widetilde W_i}{dx_i}
+
\frac{d\widetilde W_j}{dx_i}
+
\frac{d\Pi_T}{dx_i},
\qquad j\neq i.
\label{eq:wedge}
\end{equation}
The Stage-7 audit gives approximately
\[
\frac{d\widetilde W_i}{dx_i}\simeq0,
\qquad
\frac{d\widetilde W_j}{dx_i}\simeq0.50053,
\qquad
\frac{d\Pi_T}{dx_i}\simeq-0.01008,
\]
so
\[
\frac{dW^N}{dx_i}\simeq0.49045>0.
\]
Thus the repaired baseline exhibits local under-provision in the coordinated $+x_i$ direction: a marginal increase in one region's investment raises national welfare at that decentralized witness. This is not a solved first-best or global social-optimum comparison, and its sign is not implied by the sign of $\BR_i'$.

\subsection{Controlled B3--G comparison}

Table \ref{tab:welfare} reports the repaired baseline welfare comparison. B3 fixes the private price at the equilibrium fee of G; it is an identification benchmark for price endogeneity, not a literal price-regulation experiment.

\begin{table}[t]
\centering
\caption{Welfare at the repaired baseline witness}
\label{tab:welfare}
% DO NOT EDIT — GENERATED FILE
\begin{tabular}{lrrr}\toprule Environment & $W_i$ & $\Pi_T$ & $W^N$ \\ \midrule
G & 0.294879 & 0.007031 & 0.596788 \\
B3 & 0.289334 & 0.007260 & 0.585928 \\
\bottomrule\end{tabular}

\end{table}

At this repaired witness, aggregate welfare is approximately $0.59679$ in G and $0.58593$ in B3, a difference of about $0.01086$ in favor of G. This is a one-witness numerical comparison only. We do not claim that endogenous repricing generally raises welfare, that G welfare-dominates B3 over a parameter region, or that strategic substitutability is socially desirable.

\section{Robustness and Scope}\label{sec:robustness}

The robustness discussion has an expository role distinct from the proof and the repaired numerical search. Theorem \ref{thm:reversal} is a local sufficient-condition theorem on regular stationary branches, with its G branch anchored at a symmetric regular beta-zero central-interior full-game stationary state and with support-side interiority, strict routing inequalities, nonsingular local systems, and strict second-order conditions maintained locally. The repaired baseline exercise is an all-regime computational search at one parameter vector. Neither object is a general robustness theorem.

\paragraph{Network interaction.}
The analytic result organizes variation in $\beta$ near zero on continuing regular stationary branches: B3 starts from a zero cross effect with a strictly positive $\beta$ derivative, while the G branch starts from the strictly negative cross effect proved at its symmetric beta-zero anchor. The theorem establishes the existence of an implicit $\varepsilon>0$ by intersecting the neighborhoods on which branch continuation, support-side interiority, the strict routing inequalities, the matched-price path, the strict second-order conditions, and the two sign restrictions persist. It does not give a primitive formula for $\varepsilon$, prove global public optimality throughout that interval, or establish that an arbitrary parameter path remains in one allocation regime.

\paragraph{Repaired baseline numerical search.}
The repaired vector uses $\beta=0.01$, $\gamma=0.825$, and $\tau=0.35$. Stage 11R corrects the support-side welfare evaluator outside the interior region and re-runs a search over the public strategy interval, alternative routing regimes, G deviations with private repricing, B3 deviations at the fixed matched fee, multiple participation starts, the previously dangerous low-investment region, boundaries, and support-saturation neighborhoods. No profitable deviation is detected under the documented finite grid and local-refinement protocol.

This is \emph{search evidence}, not a certified global-equilibrium result. The implementation provides neither a rigorous upper bound on regret between searched points nor a proof that all possible participation or private-price continuation roots have been enumerated. It therefore does not establish a globally optimal public best response, equilibrium existence, uniqueness, a positive-measure parameter region, or a generic property.

\paragraph{Remote-public access friction.}
At the repaired vector,
\[
\kappa_L+\tau>v+\alpha.
\]
This sufficient condition makes a nonresident rival public hub dominated by non-participation for every project type and prevents the direct rival-hub free-riding deviation that invalidated the earlier draft vector. The condition is transparent but constructive. We do not claim that its magnitude is empirically calibrated or necessary for the mechanism.

\paragraph{Distribution, nonlinear network functions, and asymmetry.}
The current paper does not contain a certified theorem for arbitrary project-type distributions, arbitrary nonlinear partner-response functions, or heterogeneous regional primitives. Smoothness and strict local signs suggest that some sufficiently small perturbations may preserve the local stationary geometry, but the required function-class restrictions and global-equilibrium separation have not been independently established. We therefore treat these as extension directions rather than robustness results.

\paragraph{Numerical perturbations.}
The earlier $\pm0.5$ percent, 20-draw exercise was centered on the draft vector that later failed the all-regime deviation audit. It is non-authoritative and is not used as robustness evidence for the repaired paper. No replacement parameter-cloud exercise is needed for the current contribution because the analytic theorem and the one-vector numerical search have distinct, explicitly limited roles.

The maximum defensible scope is therefore narrow: a local analytic sign-reversal theorem under the stated regularity conditions and symmetric beta-zero G anchor, plus corrected all-regime search evidence at one baseline vector. No broader genericity, equilibrium-existence, or numerical-certification claim is made.

\section{Institutional and Empirical Bridge}
% Stage 10 substantive writing only.

\section{Related Literature}\label{sec:literature}

Our result combines ingredients that are individually familiar. The contribution is therefore best located by separating the known components from the qualitative comparison that remains.

\paragraph{Strategic commitment and downstream equilibrium responses.}
\citet{LopezVives2019} study cost-reducing R\&D with spillovers and overlapping ownership. Their benchmark model is simultaneous, while a robustness extension lets firms commit to R\&D before downstream product-market competition. That two-stage extension illustrates a broad point relevant here: a first-stage investment incentive can depend on how the downstream equilibrium responds to investment. It does not contain a shared private fee setter or two decentralized public investors, so it is not a direct predecessor of our matched-price repricing comparison. We use it only as evidence that downstream-equilibrium feedback into investment incentives is a familiar strategic mechanism.

\paragraph{Public--private competition, pricing, and investment.}
\citet{Kim2024MixedDuopoly} is a particularly close institutional and strategic benchmark: a welfare-maximizing public platform competes with a profit-maximizing private platform in a two-sided linear-city model, and the paper compares first-stage R\&D incentives followed by platform price competition. It therefore establishes directly that mixed public--private platform competition, endogenous prices, network feedback, and investment incentives can coexist in one model. Its strategic object differs from ours, however. Kim compares the investment incentives of one public and one private platform; it does not study two decentralized public investors sharing a third private alternative, nor a matched-price experiment that switches the shared private price policy off and on while holding the public--public strategic object fixed.

\citet{ZogheibBourreau2021} study competition between profit-maximizing private and welfare-oriented public firms in network infrastructure, including local public firms and investment choices. \citet{Mun2019} and \citet{MunNakagawa2010} analyze regional or cross-border infrastructure provision with government pricing and investment. More recently, \citet{LiuReshidiRivadeneyra2026} study a welfare-maximizing public payment platform competing with a profit-maximizing private platform in a two-sided market and show that the private platform can adjust fees in response to public-platform competition. These papers make clear that mixed objectives, decentralized public investment, and private price responses to public competition are not themselves novel. Relative to this line, the present paper isolates how a shared private price policy changes the sign of the strategic relationship between two public investors.

\paragraph{Two-sided platforms, network interaction, and investment.}
Strategic interaction in two-sided market oligopoly is developed in \citet{FarhiHagiu2008}. \citet{HahnKim2026} study platform investment incentives under price-parity clauses and cross-platform investment externalities, while \citet{HeEtAl2026} analyze pricing and service investment under mixed network externalities and platform competition. These contributions absorb important components of our environment: cross-side network feedback, platform pricing, and strategic investment. We do not claim novelty for those ingredients.

\paragraph{Sequential pricing and qualitative reversals.}
\citet{BontemsHamiltonLepore2025} analyze sequential pricing on multisided platforms and show that timing changes equilibrium pricing and the role of cross-platform network effects. \citet{SanchezCartas2026} shows more sharply that a follower's best response in sequential two-sided platform competition can overturn the price-skewness ranking of the simultaneous benchmark. These papers make clear that a follower response can change qualitative conclusions in platform pricing. Our narrower object is different: the follower is a shared private intermediary, the upstream strategic actors are two welfare-oriented regional governments, and the comparison holds the private fee level fixed at the full-game on-path price before switching the private price policy on. Among the closest models examined in the prior-art audit, including the mixed-platform benchmark of \citet{Kim2024MixedDuopoly}, we do not identify a result that supplies this whole matched-price public--public best-response sign comparison by direct relabeling or restriction. This positioning is intentionally narrow and does not imply a new general theory of public platforms, sequential pricing, or two-sided markets.

\section{Discussion}\label{sec:discussion}

The model isolates one mechanism: endogenous repricing by a shared private intermediary can change the local strategic relation between decentralized public investments. Several boundaries matter for interpreting that result.

First, the analytic result is local rather than global. Theorem \ref{thm:reversal} establishes the sign reversal for sufficiently small positive $\beta$ on a regular beta-zero B3 stationary branch and a regular G branch through the symmetric regular beta-zero central-interior full-game stationary state covered by the negative-cross-effect formula. The local argument also maintains support-side interiority, strict routing inequalities, nonsingular systems, the matched-price path, and strict private and public second-order conditions. It does not provide a primitive formula for the endpoint $\varepsilon$, prove an asymmetric beta-zero G result, prove global public optimality along every continuing branch, prove that such branches exist for every primitive vector, or characterize all primitives for which reversal occurs.

Second, the one-vector numerical exercise is deliberately weaker. Stage 11R corrects support-side surplus whenever participation saturates and re-runs the all-regime search with private repricing after G deviations, a fixed matched fee during B3 deviations, boundaries, low-investment histories, saturation neighborhoods, and multiple participation starts. No profitable public deviation is detected under that documented search. But the procedure does not establish a rigorous upper bound on regret over unsearched points or enumerate every possible downstream equilibrium. It is therefore search evidence, not a certified global-equilibrium or uniqueness result.

Third, the earlier draft vector is not global-equilibrium authority. It admits profitable finite deviations once routing regimes omitted by the original central-regime solver are restored. The exact rational Krawczyk calculation for that vector remains useful only as a local stationary-root diagnostic. It verifies local algebraic geometry, not global best-response optimality. The current vector is used only for the corrected numerical search and local welfare illustration.

Fourth, the novelty is proposition-level and exposed to a natural objection. Positive cross-side network feedback is familiar, negative effects transmitted through an endogenous follower price are familiar, public investments can already be strategic substitutes in fiscal-competition models, and third-party feedback can reverse strategic complementarity in other settings. The economic content here is narrower: with the private price level matched across the two environments, activating the shared private intermediary's optimizing price response can reverse the sign of the same public--public local best-response slope. We do not claim a general theory of feedback-induced strategic reversal.

Fifth, the repaired numerical vector uses a strong remote-public access friction satisfying
\[
\kappa_L+\tau>v+\alpha.
\]
Under this sufficient condition a nonresident rival public hub is dominated by non-participation for every project type. The condition prevents the direct rival-hub free-riding deviation that invalidated the draft vector. It also limits interpretation: the numerical exercise is not evidence of direct user poaching between public hubs. Public--public interaction at the reported states is mediated through the shared private route and cross-side participation. The numerical value of $\tau$ is constructive rather than empirically calibrated.

Sixth, primary-route single-homing is a modeling abstraction. Innovation projects and the firms behind them can interact with several networks. The maintained interpretation is narrower: a focal project has one principal intermediation route for the activity being modeled. Allowing project multihoming could attenuate route substitution and lies outside the present analysis.

Seventh, the private objective and instrument set are reduced-form approximations. Real commercial innovation communities may value occupancy, membership revenue, sponsorship, ecosystem growth, or longer-run strategic goals, and may use richer contracts than a single posted price. Similarly, a posted membership or access fee may bundle workspace and community services. The model uses a one-dimensional effective access price to represent the margin through which a commercial alternative responds to project demand; it does not claim robustness to arbitrary contracting environments. The sequence in which public facilitation is chosen before private pricing is likewise a maintained feature of the research question rather than a claim about all institutional timings.

Eighth, the institutional evidence establishes an analogue class, not an observed treatment-control relationship. We do not observe a causal episode in which expansion of a regional public hub induced CIC Tokyo, QWS, Level39, or another named operator to cut its price. The negative price response is a theoretical prediction. A useful empirical design would combine changes in regional facilitation capacity with project routing, partner participation, and flexible private access prices.

Finally, strategic substitution is not a policy diagnosis. The fixed-price environment is an identification benchmark, not a recommendation to regulate private prices. At the reported G stationary candidate, the corrected result pipeline finds a positive local national-welfare derivative in a coordinated $+x_i$ direction, with the own derivative computed numerically and the national derivative verified both directly and by decomposition. This is not a solved first-best problem, a global social optimum, or a general welfare theorem. No specific merger, subsidy, mandate, or price restriction is shown to be optimal.

\section{Conclusion}\label{sec:conclusion}

This paper studies two decentralized regional governments that invest in public innovation intermediation while sharing a price-setting metropolitan private alternative. Cross-side participation links project routing to partner participation. When the private access price is held fixed, this link can make the regional public investments locally strategic complements. When the private intermediary optimally reprices, its price response can dominate that positive network channel and turn the same local strategic relationship into substitution.

The main analytic result is local and sufficient. Around regular beta-zero stationary branches, with the G branch anchored at a symmetric regular beta-zero central-interior full-game stationary state and with support-side interiority, strict routing inequalities, nonsingular systems, the matched-price path, and strict second-order conditions maintained locally, the fixed-price public cross effect starts from zero and increases strictly with network interaction while the full-game cross effect is already strictly negative at that anchor state. Hence sufficiently small positive $\beta$ in the common continuation neighborhood yields opposite local best-response slopes in B3 and G. This theorem does not establish global public optimality along the continuing branches or characterize the full primitive space.

The existing repaired baseline vector is addressed separately by numerical search. Stage 11R corrects partner-surplus accounting whenever support participation saturates and re-runs the all-regime deviation search with private repricing after G deviations, a fixed matched fee during B3 deviations, boundaries, low-investment histories, support-saturation neighborhoods, and multiple participation starts. The reported stationary candidates remain near $x_G=0.83710$ and $x_{B3}=0.82589$, and no profitable public deviation is detected under the documented search protocol. Because the implementation provides no rigorous global regret bound and does not prove exhaustive enumeration of all downstream continuations, this is search evidence rather than a certified global-equilibrium, equilibrium-existence, or uniqueness result. The earlier exact rational stationary-root certificate is retained only as a local diagnostic for a draft vector that is no longer used as global-equilibrium evidence.

Welfare accounting further shows that the strategic sign is distinct from the welfare externality. The private access fee is an aggregate transfer. At the reported G stationary candidate, the corrected result pipeline computes a near-zero own reduced derivative, a positive rival-welfare derivative, and a negative private-profit derivative; their sum agrees with a direct numerical derivative of national welfare and is positive. This is a local directional statement at the reported candidate, not a solved first best or a general welfare ranking between G and B3.

The broader lesson is limited but useful. When decentralized public intermediaries share an adaptive private outside option, holding the private price level fixed can give a different qualitative picture of local public strategic interaction than allowing the shared private intermediary to reoptimize its price. How broadly this mechanism survives alternative distributions, network technologies, regional heterogeneity, contracting environments, and globally verified numerical equilibria remains an open question rather than a result of the present paper.


\section*{Acknowledgments}
OpenAI ChatGPT was used under human direction during manuscript preparation for structured drafting assistance, language editing, and code/reproducibility support; model configurations varied across the documented drafting workflow. All mathematical derivations, citations, numerical results, and final claims were independently reviewed through the reproducibility and referee-audit workflow, and the author retains full responsibility for the manuscript.

\section*{Data Availability Statement}
No external datasets were generated or analyzed in this theoretical study. The symbolic, numerical, and computer-assisted verification code underlying the reported results is included in the reproducibility materials accompanying the manuscript.

\bibliographystyle{plainnat}
\bibliography{../references/references}
\appendix
\section{Derivations and Proofs}\label{app:derivations}

\subsection{Project cutoffs}
For route $h$, write $a_{rh}=\kappa_{rh}+\mathbf 1\{h=T\}p_T$. The outside cutoff is
\[
z_{h0}^r=\frac{a_{rh}}{q_h},
\]
and the pairwise cutoff between routes $h$ and $g$, whenever $q_h\neq q_g$, is
\[
z_{hg}^r=\frac{a_{rh}-a_{rg}}{q_h-q_g}.
\]
Under the central ordering used in the main text, these expressions reduce to the cutoffs in \eqref{eq:cutoffs}.

\subsection{Private-price response}
Differentiating the private first-order condition \eqref{eq:privatefoc} with respect to $x_i$ gives \eqref{eq:px}. When $\beta=0$, private demand is linear. Let $\Delta_i$ denote the corresponding local-private quality difference and define
\[
A=\frac{1}{\Delta_1}+\frac{1}{\Delta_2},
\qquad
B=\frac{2}{q_T}.
\]
The interior private price is
\[
p_T^*
=
\frac{(\kappa_L-\kappa_T)A-\kappa_TB}{2(A+B)}.
\]
Then
\[
\frac{\partial p_T^*}{\partial A}
=
\frac{\kappa_LB}{2(A+B)^2}>0,
\qquad
A_{x_i}=-\frac{\alpha}{\Delta_i^2}<0,
\]
so $p_{T,x_i}^*<0$. This limiting case shows that endogenous private pricing does not require two-sided participation. Two-sidedness is instead essential to the positive B3 cross effect used in Theorem \ref{thm:reversal}.

\subsection{Proof of the local strategic sign reversal}\label{app:smallbeta}
Let $\delta=\alpha\beta$, $T=A_T=v+\alpha\rho_T$, $\Delta_i=A_i-A_T$, and $d=\kappa_L-\kappa_T-p_T$. On a regular interior beta-zero stationary branch, $T>0$, $\Delta_i>0$, and $d>0$. Write $M_{B3}=W_{i,x_ix_j}^{B3}$ and let $M_G$ denote the corresponding reduced full-game cross derivative.

At $\delta=0$, the B3 participation equations separate across regions once the shared private price is held fixed. Hence
\[
M_{B3}(0,x,p)=0
\]
throughout the relevant beta-zero central regime. A first-order expansion of the exact participation system and the regional objective gives
\[
\left.\frac{\partial M_{B3}}{\partial\delta}\right|_{\delta=0}
=
\frac{\alpha^2d^3}{\Delta_i^3\Delta_j^2}>0.
\]
Since $\delta=\alpha\beta$,
\[
\left.\frac{\partial M_{B3}}{\partial\beta}\right|_{\beta=0}
=
\frac{\alpha^3d^3}{\Delta_i^3\Delta_j^2}>0.
\]
Because $M_{B3}(0,x,p)$ is identically zero, differentiating along a continuing B3 stationary path and matched-price path introduces no additional first-order terms at zero: the chain-rule multipliers on the path derivatives are zero.

For G, the private price policy remains an active strategic link at $\beta=0$. At a symmetric regular beta-zero stationary state let $\Delta=A_i-A_T>0$. Substituting the exact interior private optimum into the full chain-rule decomposition \eqref{eq:repricingdecomposition} yields
\[
M_G(0)
=
-\frac{3T\alpha^2\kappa_L^2}
{16\Delta(\Delta+T)^3}<0.
\]
Thus B3 has a zero cross effect with a strictly positive derivative in $\beta$, whereas G has a strictly negative cross effect at the limiting state itself.

By the maintained nonsingularity conditions, the implicit function theorem supplies locally smooth B3 and G stationary branches and smooth reduced derivatives. Strict local second-order conditions and the central-regime inequalities also persist locally. Therefore there is an $\varepsilon>0$ such that, for every $\beta\in(0,\varepsilon)$ on the continuing stationary branches,
\[
M_{B3}(\beta)>0>M_G(\beta).
\]
Under the strict local public second-order conditions, the slope expression in \eqref{eq:brslope} has the sign of the relevant public cross derivative. Hence
\[
\BR_i^{B3\prime}>0>\BR_i^{G\prime}.
\]
This proves Theorem \ref{thm:reversal} at its stated local stationary-branch scope. The argument does not produce a primitive formula for $\varepsilon$, prove global public optimality on those branches, or establish a necessary-and-sufficient primitive characterization.

\section{Archived Exact Stationary-Root Diagnostic}\label{app:archivedexact}
An earlier draft used the primitive vector $\beta=0.05$, $\gamma=0.9$, and $\tau=0.05$. An exact rational Krawczyk calculation isolates a coupled G/B3 stationary root at that vector and certifies strict local derivative signs. Stage 4A subsequently showed that the same vector admits profitable finite public deviations once all project-routing regimes are restored. The exact calculation therefore remains mathematically valid only as a local stationary-root diagnostic. It is not evidence that the old root is a global public Nash equilibrium or subgame-perfect equilibrium, and it is not used as authority for any headline claim in the revised paper. The detailed exact intervals and verifier are retained in the repository for regression and methodological transparency.

\section{Repaired All-Regime Computational Verification}\label{app:numerics}
The repaired baseline uses $\beta=0.01$, $\gamma=0.825$, and $\tau=0.35$, with all other primitives unchanged. The independent Stage-4A evaluator reconstructs route choice from the full utility upper envelope, solves support-side participation from multiple starts, reoptimizes private pricing after public deviations, and searches the full public strategy interval $[0,1]$. It also intensively rechecks boundaries, the old low-investment failure region, and the detected maxima using additional participation starts.

The resulting symmetric witness is approximately
\[
x_G=0.83710224,
\qquad
x_{B3}=0.82589039,
\qquad
p_G=0.01840368.
\]
The best detected deviation gain is below the documented certification tolerance in each environment. Local finite-difference derivative checks at multiple step sizes give a negative G best-response slope and a positive B3 slope. The representative generated values are reported in Table \ref{tab:strategic}.

The repaired condition
\[
\kappa_L+\tau>v+\alpha
\]
analytically eliminates nonresident rival-public usage for every project type. This prevents the direct free-riding route that invalidated the old witness. The condition is sufficient for that route-dominance statement only; it is not claimed to be necessary for the sign-reversal mechanism.

No broad robustness theorem follows from this one repaired witness. In particular, the paper does not claim arbitrary-distribution robustness, arbitrary nonlinear-network robustness, generic heterogeneous-region robustness, or a positive-measure primitive region of global-equilibrium sign reversal. The earlier 20-draw perturbation exercise around the rejected vector is non-authoritative and has been removed from the paper's evidentiary chain.

\section{Additional Tables}\label{app:tables}
\begin{table}[h]
\centering
\caption{Repaired baseline parameters}
\label{tab:parameters}
% DO NOT EDIT — GENERATED FILE
\begin{tabular}{lr}\toprule Parameter & Value \\ \midrule
alpha & 0.5 \\
beta & 0.01 \\
gamma & 0.825 \\
kL & 0.27 \\
kT & 0.02 \\
rho & 0.15 \\
rhoT & 0.05 \\
tau & 0.35 \\
v & 0.1 \\
\bottomrule\end{tabular}

\end{table}

The parameter and result tables are generated from the repaired v2.1 result layer. The archived old stationary-root files are deliberately excluded from the active table pipeline.

\end{document}
